\section{Characterwise cyclotomic rank through the actual host extensions}
\label{sec:cubic-characterwise-rank-ascent}

The next reduction moves the remaining upper certificate to a host
before its large cyclic $p$-power extensions.  It does not prove that
certificate.  In particular every fixed quadratic factor and every
rational or Teichmuller coefficient character remains in the downstairs
field.  Positive class groups produced by genus theory are compatible
with the rank inequality proved here.

\paragraph{Common tower and character.}
Let $p$ be odd, let $B$ be a number field, and let $B_\infty/B$
be its cyclotomic $\mathbb Z_p$-extension.  Let $L/B$ be finite
Galois of degree prime to $p$, with group $\Delta$.  Let $T/B$
be cyclic of degree $m=p^s$, and assume
\[
                         T\cap B_\infty=B.
 \tag{G100.1}
\]
Set $N=LT$, $L_\infty=LB_\infty$ and $N_\infty=NB_\infty$.
Their finite groups are
\[
 \operatorname{Gal}(N_\infty/B_\infty)=\Delta\times P,
 \quad P=\operatorname{Gal}(T/B),\qquad
 \Gamma=\operatorname{Gal}(B_\infty/B)\simeq\mathbb Z_p.
 \tag{G100.2}
\]
Indeed $L\cap T=B$ by the coprime degrees.  Every $p$-power
quotient of $\Delta\times P$ factors through $P$, so (G100.1)
also gives $N\cap B_\infty=B$.  In particular the
$\Delta$, $P$ and cyclotomic actions commute on the fields in (G100.2).

For $J=L,N$, let $\mathcal U_J/J_\infty$ be its maximal abelian
pro-$p$ extension unramified at every finite prime and split completely
at every prime above $p$, and put
\[
 X'_J=\operatorname{Gal}(\mathcal U_J/J_\infty),\qquad
 \Omega=k[[t]],\quad t=\gamma-1,
 \tag{G100.3}
\]
where $k$ is a finite field of characteristic $p$ and $\gamma$
is a chosen topological generator of $\Gamma$.
Choose $\vartheta:\Delta\to k^\times$ and use the usual
idempotent
$e_\vartheta=|\Delta|^{-1}\sum\vartheta(\delta)^{-1}\delta$.
The modules
\[
 M_J=e_\vartheta((X'_J/p)\otimes_{\mathbb F_p}k),\qquad
 r_J=\operatorname{rank}_\Omega M_J
 \tag{G100.4}
\]
are finitely generated over $\Omega$ by the compact finite-layer
argument of AC-G74.  One may start at a common cyclotomic layer after
all $p$-primes are totally ramified.  This changes neither the nature
of the hypotheses nor the normalized-rank conclusion below.

\begin{theorem}[AC-G100-ascent: the characterwise rank does not grow faster than the host]
\label{thm:ACG100ascent}
In the preceding setting there is a natural $\Omega$-linear map
\[
                         (M_N)_P\longrightarrow M_L
 \tag{G100.5}
\]
with finite-dimensional kernel and cokernel over $k$.
It is equivariant for all the retained commuting actions before
the character projector.  Consequently
\[
                 r_L\le r_N\le m r_L.
 \tag{G100.6}
\]
The theorem permits arbitrary $r_L$; it assumes no Iwasawa
$\mu$-vanishing.  In particular it is a characterwise assertion,
not an application of a whole-field $\mu=0$ theorem.

Apply this successively to the cyclic $p$-power steps in the actual
G83 and G92 constructions, starting from a real abelian field
$F_{\rm ab}$ containing all their fixed real quadratic factors,
and the CM field $H_{\rm ab}=KF_{\rm ab}$.
Let $F_{\rm out}$ be the selected-place real host after those steps,
$H_{\rm out}=KF_{\rm out}$, and let $\alpha$ be the actual adjacent
character.  Write $r_{\rm ab},r_{\rm out}$ for its ranks as in
AC-G98, in the same cyclotomic coordinate, and
$d_{\rm ab}=[F_{\rm ab}:\mathbb Q]$,
$d_{\rm out}=[F_{\rm out}:\mathbb Q]$.
Then
\[
                 \frac{r_{\rm out}}{d_{\rm out}}
                       \le \frac{r_{\rm ab}}{d_{\rm ab}}.
 \tag{G100.7}
\]
For this application take the exact residual splitting field over
$H_{\rm ab}$, including $\chi$, $\bar\psi$ and
$\bar\epsilon^{-e}$ as needed; a larger prime-to-$p$ field
splitting these characters may also be used.

Thus a single finite downstairs strict-sector certificate
\[
 a_{{\rm ab},n}<[F_{{\rm ab},n}:\mathbb Q]
                        =d_{\rm ab}p^n
 \tag{G100.8}
\]
is sufficient for the actual G98 triangular dimension inequality
at every sufficiently high layer of the selected host.
No value of $a_{{\rm ab},n}$ is computed here.
The inequality concerns the actual character sector, not the class
number of $K(\sqrt[3]{269})$ alone.  It does not remove any fixed
quadratic factors, prove full source propagation or a prescribed
Hecke-curve containment, or supply a counting budget.
\end{theorem}

\begin{proof}
\emph{1. A compact group extension with a cyclic quotient.}
Put $\mathcal M=\mathcal U_N$, $X=X'_N$ and
$\mathscr G=\operatorname{Gal}(\mathcal M/L_\infty)$.
The defining maximality of $\mathcal M$ makes it Galois over $B$:
unramifiedness, splitting at $p$ and the abelian pro-$p$ property
are invariant under every automorphism of $N_\infty/B$.
There is an exact sequence of pro-$p$ groups
\[
                   1\longrightarrow X\longrightarrow\mathscr G
                         \longrightarrow P\longrightarrow1.
 \tag{G100.9}
\]
The action of $P$ on the abelian group $X$ is independent of lifts.
Choose a lift $\sigma\in\mathscr G$ of a generator of $P$.
Every element of $\mathscr G$ is of the form $x\sigma^i$ with
$x\in X$ and $0\le i<m$, and $\sigma^m\in X$.
All commutators are therefore generated by
$\sigma x\sigma^{-1}x^{-1}$.
Their set is the compact, hence closed, subgroup $(\sigma-1)X$.
It follows directly that
\[
        0\longrightarrow X_P\longrightarrow\mathscr G^{\rm ab}
                          \longrightarrow P\longrightarrow0.
 \tag{G100.10}
\]
This proves the needed injectivity before reduction without an
unproved group-homology vanishing assertion.
After reduction modulo $p$ the exact sequence is instead
\[
 P[p]\longrightarrow X_P/pX_P\longrightarrow
       \mathscr G^{\rm ab}/p\longrightarrow P/pP\longrightarrow0.
 \tag{G100.11}
\]
In particular the possible kernel is retained.  Its dimension is
at most one over $\mathbb F_p$, as is that of the displayed
cokernel.  Equivalently it is the usual $\operatorname{Tor}_1$
term for the finite cyclic quotient; no injectivity after mod $p$
or omission of a degree-two homology term is being assumed.

Conjugation gives canonical $\Delta\times\Gamma$ actions on
$\mathscr G^{\rm ab}$: different lifts differ by elements of
$\mathscr G$, whose inner action on its abelianization is trivial.
On $X_P$ the same independence holds because $X$ is abelian
and the $P$ action has been killed.  On $X$ itself the action of
$\Delta\times P\times\Gamma$ is canonical, since lift differences
are in $X$.  These actions commute and all maps in
(G100.10)--(G100.11) are equivariant.  No simultaneous group section
of the whole Galois extension is required.

\emph{2. The finite inertia and decomposition error.}
At a prime away from $p$, the inertia group of
$\mathcal M/L_\infty$ injects into $P$: its kernel is inertia
for $\mathcal M/N_\infty$, which is trivial.
Only primes ramified in the fixed finite extension $N/L$ can
contribute.  There are finitely many primes above this set in
$L_\infty$.  Indeed every finite prime is finitely decomposed
in a cyclotomic $\mathbb Z_p$-extension; away from $p$ its
Frobenius image is nonzero because a positive integer norm larger
than one cannot be a root of unity in $\mathbb Z_p^\times$.
The primes above $p$ are also finitely many, and eventually totally
ramified in the cyclotomic tower.

At a prime above $p$, the full decomposition group of
$\mathcal M/L_\infty$ likewise injects into $P$, because
$\mathcal M/N_\infty$ is split there.  This uses full
$p$-decomposition groups, not merely their inertia.
In $\mathscr G^{\rm ab}$ the image of each such inertia or
decomposition group is independent of the chosen prime above it.
Let $R$ be the subgroup generated by all these images, for the
finitely many relevant base primes.  Each image is a finite
$p$-group of exponent dividing $m$; the ambient group is abelian.
Their sum is therefore finite and closed.
It is stable under $\Delta\times\Gamma$.

There is an exact identification
\[
                 \mathscr G^{\rm ab}/R\simeq X'_L.
 \tag{G100.12}
\]
One direction follows because this quotient is abelian pro-$p$,
unramified away from $p$, and trivial on the full decomposition
groups above $p$.  Conversely, any abelian pro-$p$ extension
of $L_\infty$ unramified everywhere and split at $p$ remains
so after base change to $N_\infty$.  That compositum is contained
in $\mathcal M$ by its defining maximality.  The original extension
is therefore already contained in $\mathcal M$.
Thus the quotient is the full maximal extension defining $X'_L$,
not merely an unspecified quotient of its class module.
This also proves equivariance of (G100.12).
Reduction gives a surjection
\[
       \mathscr G^{\rm ab}/p\longrightarrow X'_L/p
                  \quad\hbox{with kernel a quotient of }R/pR.
 \tag{G100.13}
\]
Combining it with (G100.11) gives a map
$X_P/p\to X'_L/p$ with finite kernel and cokernel.
The finite error modules have continuous cyclotomic action.
After tensoring with $k$, a finite quotient of the pro-$p$ group
$\Gamma$ acts, so some power of $t=\gamma-1$ kills them.
Taking $e_\vartheta$ is exact since $|\Delta|$ is prime to $p$.
Also $P$-coinvariants commute with reduction, scalar extension
and this projector.  This proves (G100.5) and, on putting
$Q=k((t))$, the actual isomorphism
\[
                   (M_N\otimes_\Omega Q)_P
                               \simeq M_L\otimes_\Omega Q.
 \tag{G100.14}
\]

\emph{3. Cyclic $p$-group Nakayama and the rank inequality.}
Let $V=M_N\otimes_\Omega Q$.  This is finite-dimensional over
$Q$ and has a commuting $P$ action.  The augmentation ideal of
$Q[P]$ is generated by $\sigma-1$, and
\[
                         (\sigma-1)^m=\sigma^m-1=0
 \tag{G100.15}
\]
in characteristic $p$.
Choose $r_L$ vectors lifting a basis of $V_P$, whose dimension
is $r_L$ by (G100.14).  Their $Q[P]$-span has quotient $W$
satisfying $W=(\sigma-1)W$.  Iterating (G100.15) gives $W=0$.
Thus there is a surjection $Q[P]^{r_L}\to V$, whence
$\dim_QV\le m r_L$.
The quotient map $V\to V_P$ gives the other inequality.
This proves (G100.6) without a freeness or semisimplicity
assumption on the characteristic-$p$ $P$-module.

\emph{4. Apply the reduction before the actual host construction.}
In G83, collect all of its fixed real quadratic factors into
$F_{\rm ab}$.  In its positive-sign branch this includes the
real subfield of the G81 multiquadratic host; in the odd branch
it also includes the fresh real quadratic used in the Kummer
prime prescription.  In both branches it contains the real
quadratic base $E$ of the ray-class extension.
The field $F_{\rm ab}/\mathbb Q$ is real abelian.  The
G83 extension $T_E/E$ is cyclic of $p$-power degree, so
$F_{\rm ab}T_E/F_{\rm ab}$, and its base change to
$H_{\rm ab}$, are cyclic $p$-power extensions.
They are unramified at $p$.  The initial coefficient-prime
fields are unramified over $\mathbb Q_p$, whereas their
cyclotomic towers are totally ramified there.
Consequently these finite extensions meet the cyclotomic tower
trivially, verifying (G100.1).

The G92 extension over the resulting real host is cyclic of
$p$-power degree and splits at every prime above $p$.
It too meets the cyclotomic tower trivially: a common subextension
would be both split and totally ramified at those primes.
Its CM base change has the same property.
At each of these steps the residual splitting extension has
prime-to-$p$ degree, hence is disjoint from the cyclic host
extension.  The Galois group is the product $\Delta\times P$
used in the proof.  Thus no passage to a possibly noncommuting
normal closure is needed for this application.
Multiplying (G100.6) along these steps gives (G100.7), since
the real degree $d$ is multiplied by the same factors.

If the tower is reindexed by replacing its base by the $p^n$th
layer, its new parameter is
$t_n=(1+t)^{p^n}-1=t^{p^n}$ in characteristic $p$.
Restriction of scalars multiplies $\Omega$-rank by $p^n$,
while the real degree is multiplied by $p^n$.
Finite torsion modules remain torsion.  Hence the normalized rank
is unchanged by the finite cyclotomic layers used in G91/G93-fill.

Finally the actual adjacent character is nontrivial on every
coefficient-prime inertia group already at $H_{\rm ab}$:
its cubic tame part has order three, and the rational part has
inertial order dividing $268$.  Their orders cannot cancel.
The fixed quadratic factors are unramified at the coefficient
prime; the later local extensions have $p$-power degree or are
unramified.  They do not kill the order-three inertia character.
Thus the permutation correction in AC-G74 vanishes downstairs
as well as upstairs, exactly as in AC-G98.  Its finite-layer
rank certificate gives
\[
                         r_{\rm ab}p^n\le a_{{\rm ab},n}.
 \tag{G100.16}
\]
Equation (G100.8) implies $r_{\rm ab}<d_{\rm ab}$; (G100.7)
then gives $r_{\rm out}<d_{\rm out}$.
The rational central and trivial sectors on the selected host
are bounded by AC-G98.  Its theorem therefore gives the stated
eventual triangular inequality.  The genus-positive adjacent and
central classes of G93 are not lost: the later cyclotomic
restriction maps for their nontrivial prime-to-$p$ characters
are injective, as already proved in G93-fill.
No upper estimate was obtained by using that positivity.
\end{proof}

\paragraph{Exact source and remaining scope.}
The maximal unramified $p$-split extension, finite-layer class map,
compact finite generation and character orientation are those of
AC-G74, whose complete applied proof is retained.  The finite
cyclic group extension, its abelianization, all error groups and
the localized rank inequality are proved here.
For finite decomposition and eventual total ramification in the
cyclotomic tower, the exact primary input is Sharifi,
\emph{Iwasawa Theory}, Propositions~3.1.9 and~3.1.11, at
\url{https://www.math.ucla.edu/~sharifi/notes/iwasawa-ch03.html}.
Their complete applied proofs were freshly reread for this theorem,
in addition to the earlier G74 source check.  Base change
of unramified and split extensions is the usual valuation/class-field
property, within the already retained global reciprocity interface.
No Iwasawa Riemann--Hurwitz formula or new $\mu$-vanishing theorem
is used.

The fixed-field class-number certificate in AC-K31 concerns
$K(\sqrt[3]{269})$ and its untwisted pro-$269$ tower.
It does not compute the $\alpha$-sector after the necessary
prime-to-$269$ characters and all fixed quadratic factors are
adjoined.  Those are exactly the remaining inputs to (G100.8).
The result removes the cyclic $p$-host enlargement from that
upper-certificate problem.  It does not establish the certificate,
the whole changed-tuple propagation theorem, a prescribed
varying-character Hecke curve, or the original counting budget.

\paragraph{Geometric dictionary.}
The field $\mathcal U_J$ classifies abelian pro-$p$ covers of the
arithmetic curve unramified at finite places and split over its
coefficient-prime points.  Passing through the cyclic host cover
and taking coinvariants introduces only the finite inertia and
full-decomposition images explicitly killed in (G100.12).
This uses the already audited SGA~1 finite-cover dictionary and
asserts no new archival theorem.
